\section{Business insights}

\subsection{Insight 1: Stockout làm sai lệch tín hiệu demand}

Tỷ lệ stockout row khoảng 24.65\% cho thấy đây không phải hiện tượng hiếm. Nếu doanh nghiệp chỉ quan sát sales, một phần demand bị ẩn sẽ bị hiểu nhầm là demand thấp. Trong vận hành tồn kho, đây là vấn đề nghiêm trọng vì forecast thấp dẫn tới reorder thấp.

\subsection{Insight 2: Latent demand recovery cho biết mức demand bị bỏ lỡ}

Recovered demand tăng khoảng 10.68\% so với observed sales. Con số này không nên hiểu là doanh nghiệp chắc chắn mất đúng từng đó doanh số, vì recovered demand là proxy. Tuy nhiên, nó là chỉ báo hữu ích: observed sales có khả năng đang understated demand ở mức đáng kể.

\subsection{Insight 3: Weekly seasonality là baseline nghiệp vụ mạnh}

Seasonal naive 7-day rất mạnh. Về business, điều này nghĩa là nhu cầu trong bán lẻ có lịch tuần rõ ràng. Khi triển khai hệ thống forecast, không nên bỏ qua các quy tắc đơn giản như lấy demand cùng kỳ tuần trước. ML chỉ có ý nghĩa khi nó cải thiện trên nền seasonality mạnh đó hoặc kết hợp tốt với nó.

\subsection{Insight 4: Model cuối nên hỗ trợ replenishment}

Mục tiêu cuối không phải tạo model có metric đẹp nhất trong phòng thí nghiệm, mà là hỗ trợ quyết định nhập hàng. Forecast recovered demand phù hợp hơn observed sales vì nó hướng tới nhu cầu vận hành. Hybrid seasonal-ML có thể xem là một forecast thực dụng: dùng seasonality làm nền, dùng ML để hiệu chỉnh theo stockout và pattern gần đây.

\subsection{Khuyến nghị triển khai}

Nếu triển khai thực tế, hệ thống nên báo cáo song song observed sales forecast và recovered demand forecast. Khi hai giá trị lệch nhau lớn ở các SKU có stockout cao, đây là tín hiệu cần kiểm tra replenishment. Doanh nghiệp cũng nên theo dõi uplift do imputation theo SKU/store để phát hiện nhóm sản phẩm có rủi ro mất doanh số do hết hàng.
